% GENERATED FILE. Edit canonical lessons, then run npm run generate:books.
% canonical-source-hash: fnv1a64:b49ec5713207aca2
% canonical-lessons: IT-C469-bufera, IT-C469-uragano, IT-C469-raffica, IT-C469-soffiare, IT-C469-nubifragio

\chapter{Vento e tempesta}
\label{ch:it-vento-e-tempesta}
\hlchaptermodality{469}

\begin{chapteropening}
\textbf{By the end of this chapter:} \emph{I can describe wild weather: a blizzard, a hurricane, a gust of wind, the wind blowing and a cloudburst.}

Complete the last lesson of chapter 469: I can describe wild weather: a blizzard, a hurricane, a gust of wind, the wind blowing and a cloudburst.
\end{chapteropening}

\section[la bufera]{la bufera — the blizzard, the storm}

\label{lesson:IT-C469-bufera}



\begin{quote}
Before the new one: say the Italian for the dew, then the Italian for to deep-freeze.
\end{quote}

\subsection*{You'll want to know: la bufera}
\textbf{la bufera} — \textquotedblleft{}the blizzard, the storm\textquotedblright{}.

\begin{quote}
In montagna c'è stata una bufera di neve.
\end{quote}

\begin{quote}
\emph{In the mountains there was a snowstorm.}
\end{quote}

\emph{Una bufera di vento} — a gale.

\begin{cousinweb}
\emph{bufera} is built on \emph{bufare} 'to blow' (EtymDB), a word that imitates a gust of wind; the origin is not settled beyond that.
\end{cousinweb}

\begin{grammarlens}[title={What you've built}]
One more word for telling, explaining and describing at B1.
\end{grammarlens}

\subsection*{Your turn}
\begin{itemize}
\raggedright
  \item \emph{Say it:} \emph{la bufera}
  \item \emph{Say it:} \emph{la bufera}, once more
  \item \emph{Say it:} a sentence of your own with \emph{la bufera}
  \item \emph{Recall:} say the Italian for the dew, then the Italian for to deep-freeze, then say \emph{la bufera} again
\end{itemize}

\begin{tcolorbox}[breakable,colback=teal!4,colframe=teal!35!black,title={Before you move on}]
What does la bufera mean? (\textquotedblleft{}The blizzard, the storm\textquotedblright{}.) Say it once more.
\end{tcolorbox}
\section[l'uragano]{l'uragano — the hurricane}

\label{lesson:IT-C469-uragano}



\begin{quote}
Before the new one: say the Italian for to deep-freeze, then the Italian for the blizzard.
\end{quote}

\subsection*{You'll want to know: l'uragano}
\textbf{l'uragano} — \textquotedblleft{}the hurricane\textquotedblright{}.

\begin{quote}
L'uragano ha distrutto molte case.
\end{quote}

\begin{quote}
\emph{The hurricane destroyed many houses.}
\end{quote}

Also a person who sweeps in like one.

\begin{cousinweb}
\emph{uragano} was borrowed from Spanish \emph{huracán} (EtymDB), which took it from a language of the Caribbean.
\end{cousinweb}

\begin{grammarlens}[title={What you've built}]
One more word for telling, explaining and describing at B1.
\end{grammarlens}

\subsection*{Your turn}
\begin{itemize}
\raggedright
  \item \emph{Say it:} \emph{l'uragano}
  \item \emph{Say it:} \emph{l'uragano}, once more
  \item \emph{Say it:} a sentence of your own with \emph{l'uragano}
  \item \emph{Recall:} say the Italian for to deep-freeze, then the Italian for the blizzard, then say \emph{l'uragano} again
\end{itemize}

\begin{tcolorbox}[breakable,colback=teal!4,colframe=teal!35!black,title={Before you move on}]
What does l'uragano mean? (\textquotedblleft{}The hurricane\textquotedblright{}.) Say it once more.
\end{tcolorbox}
\section[la raffica]{la raffica — the gust (of wind); the burst}

\label{lesson:IT-C469-raffica}



\begin{quote}
Before the new one: say the Italian for the blizzard, then the Italian for the hurricane.
\end{quote}

\subsection*{You'll want to know: la raffica}
\textbf{la raffica} — \textquotedblleft{}the gust (of wind); the burst\textquotedblright{}.

\begin{quote}
Una raffica di vento mi ha portato via il cappello.
\end{quote}

\begin{quote}
\emph{A gust of wind blew my hat away.}
\end{quote}

\emph{Una raffica di domande} — a barrage of questions.

\begin{cousinweb}
The origin of \emph{raffica} is not settled.
\end{cousinweb}

\begin{grammarlens}[title={What you've built}]
One more word for telling, explaining and describing at B1.
\end{grammarlens}

\subsection*{Your turn}
\begin{itemize}
\raggedright
  \item \emph{Say it:} \emph{la raffica}
  \item \emph{Say it:} \emph{la raffica}, once more
  \item \emph{Say it:} a sentence of your own with \emph{la raffica}
  \item \emph{Recall:} say the Italian for the blizzard, then the Italian for the hurricane, then say \emph{la raffica} again
\end{itemize}

\begin{tcolorbox}[breakable,colback=teal!4,colframe=teal!35!black,title={Before you move on}]
What does la raffica mean? (\textquotedblleft{}The gust (of wind); the burst\textquotedblright{}.) Say it once more.
\end{tcolorbox}
\section[soffiare]{soffiare — to blow}

\label{lesson:IT-C469-soffiare}



\begin{quote}
Before the new one: say the Italian for the hurricane, then the Italian for the gust.
\end{quote}

\subsection*{You'll want to know: soffiare}
\textbf{soffiare} — \textquotedblleft{}to blow\textquotedblright{}.

\begin{quote}
Il vento soffia forte dal nord.
\end{quote}

\begin{quote}
\emph{The wind is blowing hard from the north.}
\end{quote}

\emph{Soffiarsi il naso} — to blow your nose.

\begin{cousinweb}
\emph{soffiare} is inherited from Latin \emph{sufflāre} 'to blow into' (EtymDB), from \emph{flāre} 'to blow'.
\end{cousinweb}

\begin{grammarlens}[title={What you've built}]
One more word for telling, explaining and describing at B1.
\end{grammarlens}

\subsection*{Your turn}
\begin{itemize}
\raggedright
  \item \emph{Say it:} \emph{soffiare}
  \item \emph{Say it:} \emph{soffiare}, once more
  \item \emph{Say it:} a sentence of your own with \emph{soffiare}
  \item \emph{Recall:} say the Italian for the hurricane, then the Italian for the gust, then say \emph{soffiare} again
\end{itemize}

\begin{tcolorbox}[breakable,colback=teal!4,colframe=teal!35!black,title={Before you move on}]
What does soffiare mean? (\textquotedblleft{}To blow\textquotedblright{}.) Say it once more.
\end{tcolorbox}
\section[il nubifragio]{il nubifragio — the cloudburst, the downpour}

\label{lesson:IT-C469-nubifragio}



\begin{quote}
Before the new one: say the Italian for the gust, then the Italian for to blow.
\end{quote}

\subsection*{You'll want to know: il nubifragio}
\textbf{il nubifragio} — \textquotedblleft{}the cloudburst, the downpour\textquotedblright{}.

\begin{quote}
Il nubifragio ha allagato le strade.
\end{quote}

\begin{quote}
\emph{The cloudburst flooded the streets.}
\end{quote}

A lot of rain in a short time.

\begin{cousinweb}
\emph{nubifragio} is a learned borrowing from Medieval Latin \emph{nūbifragium} (EtymDB), from \emph{nūbēs} 'cloud' and \emph{frangere} 'to break': the cloud breaks.
\end{cousinweb}

\begin{grammarlens}[title={What you've built}]
One more word for telling, explaining and describing at B1.
\end{grammarlens}

\subsection*{Your turn}
\begin{itemize}
\raggedright
  \item \emph{Say it:} \emph{il nubifragio}
  \item \emph{Say it:} \emph{il nubifragio}, once more
  \item \emph{Say it:} a sentence of your own with \emph{il nubifragio}
  \item \emph{Recall:} say the Italian for the gust, then the Italian for to blow, then say \emph{il nubifragio} again
\end{itemize}

\begin{tcolorbox}[breakable,colback=teal!4,colframe=teal!35!black,title={Before you move on}]
What does il nubifragio mean? (\textquotedblleft{}The cloudburst, the downpour\textquotedblright{}.) Say it once more.
\end{tcolorbox}
